\chapter{A Margin for Force Closure}
\label{chap:margin}

% TODO: This is the core theoretical contribution chapter.
% Suggested structure:
%  - Motivate why a binary force-closure test is insufficient for
%    path planning (no notion of "how robust" a configuration is).
%  - Survey candidate margin/quality metrics from the grasping
%    literature (e.g. largest inscribed ball in the wrench space,
%    minimum singular value formulations, etc.) and state your
%    selection criteria (computational cost, differentiability,
%    physical interpretability, suitability as a planning cost).
%  - Derive/define your chosen force closure margin precisely:
%    inputs, formula, units, bounds, degenerate cases.
%  - Prove or argue its key properties (continuity, monotonicity
%    with respect to contact quality, behaviour at the force
%    closure boundary, computational complexity).
%  - Discuss how the margin will be consumed by a path planner
%    (e.g. as a cost term, as a constraint, as a heuristic).

\section{Requirements for a Path-Planning-Ready Metric}

\section{Candidate Metrics}
Potential: \textbf{[The-Synthesis-of-Stable-Force-Closure-Grasps]}

\textbf{[Handbook-of-Robotics]} Example 3 shows hyperstatical graps, is that relevant for the snake case as well? So LP2 valid, but not LP3. Need the snake JAcobian, so joint and angle configurations to calculate LP3. 

\section{Definition of the Force Closure Margin}
\subsection{Wrench Space and Force Closure}
\label{sec:wrench-space}

Consider a planar contact configuration between the snake robot's body and
$n_c$ obstacle contact points $\mathbf{p}_i \in \mathbb{R}^2$, $i = 1,
\dots, n_c$, with unit outward normals $\mathbf{n}_i$ and tangents
$\mathbf{t}_i = \mathbf{n}_i^{\perp}$. Each contact is modeled as a
point contact with Coulomb friction of coefficient $\mu$, transmitting a
planar force
\begin{equation}
\mathbf{f}_i = f_{n,i}\,\mathbf{n}_i + f_{t,i}\,\mathbf{t}_i,
\qquad f_{n,i} \ge 0,\ \ |f_{t,i}| \le \mu f_{n,i},
\label{eq:friction-cone}
\end{equation}
but no moment about the contact point (hard-finger contact model).

Every admissible contact force induces a wrench $(f_x, f_y, \tau)$ on the
robot body about a reference point $P_0$ (e.g.\ the body's center of mass),
with $\tau = (\mathbf{p}_i - P_0) \times \mathbf{f}_i$. Collecting the
normal- and tangential-force contributions of all contacts into the
\emph{grasp matrix} $G \in \mathbb{R}^{3 \times 2n_c}$
(Sec.~\ref{sec:grasp-matrix}) gives the linear map from contact force
coefficients $\mathbf{c} = [f_{n,1}, f_{t,1}, \dots, f_{n,n_c},
f_{t,n_c}]^\top$ to the resulting body wrench:
\begin{equation}
\mathbf{w} = G\,\mathbf{c}, \qquad \mathbf{w} = (f_x, f_y, \tau) \in \mathbb{R}^3.
\label{eq:wrench-map}
\end{equation}

Each column of $G$ is therefore a \emph{primitive wrench} — the wrench
produced by a unit force along one basis direction (normal or tangential)
at one contact. The set of all wrenches the contact configuration can
actually produce, once the friction cone constraints~\eqref{eq:friction-cone}
are imposed on $\mathbf{c}$, is the \emph{wrench space} (or wrench cone)
of the grasp:
\begin{equation}
\mathcal{W} = \left\{\, G\mathbf{c} \ \middle|\ \mathbf{c} \text{ satisfies}~\eqref{eq:friction-cone} \text{ for all } i \,\right\} \subset \mathbb{R}^3.
\label{eq:wrench-space}
\end{equation}
Because the friction-cone constraints are conic (positively homogeneous),
$\mathcal{W}$ is a convex cone spanned by the primitive wrenches, obtained
as the Minkowski sum of the (linearized) friction cones mapped through $G$.

\subsection{LP2}
\label{sec:force-closure-margin}

Following the frictional form-closure test of
Bicchi~et~al. cite{bicchi2008grasping} (Springer Handbook of Robotics,
Sec.~38.4.2), the friction cone~\eqref{eq:friction-cone} at each contact
is linearized into a set of half-space constraints $F\mathbf{c} \le
\mathbf{0}$, where $F$ collects the facet normals of every contact's
linearized friction cone. Force closure — i.e.\ $\mathbf{0} \in
\operatorname{int}(\mathcal{W})$ — is then tested constructively by
searching for a nonzero $\mathbf{c}$ that produces zero net wrench while
sitting strictly inside every friction cone, via the linear program

\begin{equation}
\begin{aligned}
\underset{\mathbf{c},\, d}{\text{maximize}} \quad & d \\
\text{subject to} \quad & G\mathbf{c} = \mathbf{0}, \\
& F\mathbf{c} \le -\mathbf{1}\, d, \\
& d \ge 0, \\
& \mathbf{e}^\top \mathbf{c} \le n_c,
\end{aligned}
\label{eq:lp2}
\end{equation}

where $\mathbf{1}$ and $\mathbf{e}$ are vectors of ones of appropriate
dimension, and the normalization constraint $\mathbf{e}^\top \mathbf{c}
\le n_c$ prevents the trivial unbounded scaling of $\mathbf{c}$.

The optimal value $d^\star$ of~\eqref{eq:lp2} is taken as the
\emph{force closure margin}:

\begin{equation}
d^\star
\begin{cases}
= 0, & \text{no frictional force closure } (\mathbf{0} \notin \operatorname{int}(\mathcal{W})), \\
> 0, & \text{frictional force closure holds } (\mathbf{0} \in \operatorname{int}(\mathcal{W})),
\end{cases}
\label{eq:fc-margin-interpretation}
\end{equation}

with larger $d^\star$ indicating that the origin sits further inside the
wrench space $\mathcal{W}$ — i.e.\ a more robust closure, better able to
withstand disturbances or produce the propulsive wrench required by the
locomotion gait. This margin serves as the continuous force-closure
metric used in the path planner,
candidate configurations are scored (or filtered) by $d^\star$, computed
once per contact set by solving the linear program~\eqref{eq:lp2}.

\section{Properties of the snake robot}
Should follow the defininition in Gravdahl 2025.
As in [Gravdahl, Nordal] the link width is neglected; for frictionless contacts this is exact, while with friction it slightly changes the moment arms of the tangential contact forces.

\section{Properties of the Margin}

\section{Using the Margin in Path Planning}

